% Created 2025-06-04 Wed 23:44
% Intended LaTeX compiler: pdflatex
\documentclass[11pt]{article}
\usepackage[utf8]{inputenc}
\usepackage[T1]{fontenc}
\usepackage{graphicx}
\usepackage{longtable}
\usepackage{wrapfig}
\usepackage{rotating}
\usepackage[normalem]{ulem}
\usepackage{amsmath}
\usepackage{amssymb}
\usepackage{capt-of}
\usepackage{hyperref}
\author{Juan Ignacio Raggio}
\date{\today}
\title{Errores comunes segundo parcial poo}
\hypersetup{
  pdfauthor={Juan Ignacio Raggio},
  pdftitle={Errores comunes segundo parcial poo},
  pdfkeywords={},
  pdfsubject={},
  pdfcreator={Emacs 29.4 (Org mode 9.7.29)}, 
  pdflang={English}}
\begin{document}

\maketitle
\tableofcontents

\section{Generales}
\label{sec:org76c486c}
\begin{itemize}
  \item \textbf{No olvidarse de implementar metodos}
    toString, to\textsubscript{s}, etc.
  \item Leer muy detenidamente el test y no olvidarse metodos que pidan implementar
\end{itemize}
\subsection{<=>(other) vs ==(other) vs equals(other) - Comparable}
\label{sec:org30518ee}
\begin{itemize}
  \item Ojo porque en caso de ruby si a <=> se le pasa un objeto que no es instancia de self, retorna:
    \textbf{nil}

  \item Si a ==(other) o a equals(other) le pasas una no-instancia-self
    \textbf{false}
\end{itemize}
\subsubsection{Conclusion de pasaje de no instancia:}
\label{sec:orge7a800b}
\begin{center}
  \begin{tabular}{lll}
    Tipo de metodo & Ruby & Java\\
    Comparacion & nil & No compila\\
    Equidad & false & false\\
  \end{tabular}
\end{center}
\section{Java}
\label{sec:org818f731}
\subsection{Prestarle mucha atencion a la coleccion que se devuelve}
\label{sec:org9fc9d4b}
Es la parte facil pero al hacerla rapido se puede meter la pata
\begin{itemize}
  \item \uline{Retornar una Coleccion / Mapa} que cumpla ciertas condiciones
  \item Acordarse que \[Map\ \neq\ Collection\ \land\ Map\ \nsubseteq\ Iterable\]
  \item Por lo tanto es muy importante que si tengo que devolver un mapa, devuelvo un mapa, nisiquiera un Iterable porque no es Iterable un Mapa:
\end{itemize}
\section{Ruby}
\label{sec:orgdf982cf}
\subsection{Leer bien los test y \underline{Los codigos que te pueden dar} para que sean parte de tu implementacion}
\label{sec:orgc2c4d0e}
\begin{itemize}
  \item A veces aparecen clases que tenes que implementar y no son tan explicitas (como en el de passHolder)
  \item El metodo map puede ser muy util:
\end{itemize}

\begin{verbatim}


class LimitedDaysPassholders < ThemeParkPassholders
 def initialize(visitor, days_cap)
   super(visitor)
   @days_cap = days_cap
 end


 def add_pass?(park, day)
   # Donde @passes es un SortedSet
   # y day es un atributo de los objetos que se guardan @passes
   super && @passes.length < @days_cap &&
   !@passes.map{ |pass| pass.day }.include?(day)
 end
end

\end{verbatim}
\subsection{to\textsubscript{s}}
\label{sec:org196c29d}
\begin{itemize}
  \item No olvidarse de hacer los to\textsubscript{s}
  \item Es mejor hacer interpolacion de strings que concatenacion
\end{itemize}
\begin{verbatim}
i = "interpolacion"
"Esto es una #{i} de strings"
\end{verbatim}
\subsection{Metodos protected}
\label{sec:orgf2f23bc}
\begin{itemize}
  \item En caso de que un metodo necesite leer una variable de instancia de otra instancia de la misma clase:
    Es \textbf{MUY IMPORTANTE} hacer que esos getters sean \texttt{protected} porque sino estariamos leakeando variables sin sentido ( ojo, puede que en algun caso se quiera dejar normal )
\end{itemize}
\subsection{Clases abstractas}
\label{sec:org72e9d32}
Sabemos de objetos que las clases abstractas \texttt{NUNCA} pueden instanciarse, por lo que el metodo \texttt{initialize} de ruby en una ``clase asbtracta'' a pesar de que no existan formalmente, los vamos a simular haciendo que el initialize tenga siempre la siguiente forma:

\begin{verbatim}
class pirlas
  def initialize
    raise 'Este metodo debe sobreescribirse'
  end

  Otros metodos del pirlas
end
\end{verbatim}
\subsubsection{{\bfseries\sffamily DONE} DOUBT Aca siempre me tira warnings el ide, entiendo porque es, no entiendo entonces porque estamos forzando a hacerlo de esta manera}
\label{sec:org4a7a5e7}
Si, no importan los warnings para clases abstractas
\subsubsection{{\bfseries\sffamily DONE} DOUBT En el ejercicio de Accounts del tp8 no se porque crean un metodo de la siguiente forma:}
\label{sec:org742a6ee}
\begin{verbatim}
def extract?(_amount)
  raise '...'
end
\end{verbatim}
\begin{itemize}
  \item[{$\boxtimes$}] Porque el \_? y porque cuando ``sobreescriben'' este metodo no lo ponen? es para que no tire warnings?
    El \_ se usa cuando ese parametro no lo vas a usar dentro del metodo
  \item[{$\boxtimes$}] En el initialize igualmente seguiriamos con el mismo problema
    Si, no importan los warnings para clases abstractas
\end{itemize}
\subsection{Colecciones}
\label{sec:org7ab8b88}
\uline{Muy importante}:
\begin{itemize}
  \item Es muy comun cuando quiero extraer info de las colecciones, ya sea: Hash, Set, etc. Tener que \textbf{pasarlas a array}, por ejemplo:
    En este caso, tenemos @bag que es un Hash table y cuando queremos obtener las llaves para ver si estamos en el final, usamos .keys para obtenerlas
\end{itemize}
\begin{verbatim}
def to_s
  keys = @bag.keys
  s = ''
  keys.each {
    s += "#{it.to_s} => #{@bag[it].to_s}"
    s += ", " unless it == keys.last
  }
  "{#{s}}"
end
\end{verbatim}
\subsubsection{No olvidarse de sobreescribir ==(other) y hash}
\label{sec:orge293cd6}
\textbf{Retornan false si other no es instancia}
\begin{itemize}
  \item \textbf{OJO} porque cuando solo se sobreescribe al ==(other) \uline{NO OLVIDARSE DEL ALIAS}
    Como las colecciones usan eql? para comparar, si no lo sobreescribimos no funciona
\end{itemize}
\begin{verbatim}
alias_method :eql?, :==
\end{verbatim}
\end{document}
